\section{Introduction and statement of results}
Painlev\'e metrics \cite{Pain1897, Per1990} are a class of Riemannian metrics that provide a broad generalization of the well-known separable metrics of St\"ackel \cite{Eis1934, Sta1897} to the case in which the Hamilton--Jacobi equation for the geodesic flow can be additively separated into partial differential equations depending on groups of independent variables rather ordinary differential equations resulting from a complete orthogonal separation. In particular, while St\"ackel metrics in dimension $n$ admit $n$ linearly independent Poisson-commuting quadratic first integrals of the geodesic flow, Painlev\'e metrics in dimension $n$ will admit only $r<n$ such integrals in general, where $r$ denotes the number of groups of variables.

Our goal in this paper is to carry out the extension to Painlev\'e metrics of the well-known results \cite{Ben2016, BCR1-2002,BCR2-2002, Eis1934, KaMi1980, KaMi1981, KaMi1984} which relate in the St\"ackel case the additive separation of variables for the Hamilton--Jacobi equation to the multiplicative separation of variables for the Helmholtz equation, and the existence of quadratic first integrals of the geodesic flow to that of symmetry operators for the Laplace--Beltrami operator. We shall thus obtain the generalization to Painlev\'e metrics of the Robertson separability conditions~\cite{Rob1928} for the Helmholtz equation for St\"ackel metrics, and a formulation of these generalized Robertson conditions in terms of the vanishing of the off-block diagonal components of the Ricci tensor, thereby extending the classical result proved by Eisenhart \cite{Eis1934} for St\"ackel metrics. We shall also show that when the generalized Robertson conditions are satisfied, there exist $r<n$ linearly independent second-order differential operators which commute between themselves and with the Laplace--Beltrami operator. These operators will be shown to admit the block-separable solutions of the Helmholtz equation as formal eigenfunctions, with the separation constants arising from the separation into groups variables as eigenvalues. Finally, we shall also discuss conformal deformations of Painlev\'e metrics satisfying a further generalization of the Robertson conditions, which is compatible with the separation into blocks of variables of the Helmholtz equation, leading to solutions which are $R$-separable in blocks.

Before describing our results in further detail, we should remark that independently of the interest of Painlev\'e metrics from the point of view of separation of variables, a key motivation for our study lies in the goal of constructing geometric models of manifolds with boundary endowed with Painlev\'e metrics, with the goal of investigating the anisotropic Calder\'on problem in this class of geometries. Recall that the anisotropic Calder\'on problem consists in recovering the metric of a Riemannian manifold with boundary from the knowledge of the Dirichlet-to-Neumann map defined by the Laplace--Beltrami operator. The anisotropic Calder\'on problem is at the center of a great amount of current research activity. We refer to \cite{DKN2016,DKN2019a, GT2013, KKL2001, KS2014, LU2001, LeU1989, Sa2013, Uhl2009, Uhl2014} and the references therein for surveys of recent results on this problem. We have recently investigated the anisotropic Calder\'on problem at fixed energy for geometric models consisting of classes of St\"ackel manifolds with boundary, where the separation of variables for the Helmholtz equation allows the decomposition of the Dirichlet-to-Neumann map onto a~basis of joint eigenfunctions of the symmetry operators resulting from the complete separation of variables, enabling us to obtain a series of uniqueness and non-uniqueness results for the Calder\'on problem, with no a-priori assumptions of analyticity, or on the existence of isometries \cite{DKN2015, DKN2016, DKN2017, DKN2018a, DKN2019b, Gob2016}. In the case of Painlev\'e metrics, the separation of the Helmholtz into groups of variables and the concomitant families of commuting symmetry operators admitted by these metrics will serve as an effective starting point for the investigation of the anisotropic Calder\'on problem in this more general setting.

In order to put the results of the present paper in context, we first briefly recall some well-known definitions and results pertaining to St\"ackel metrics and their separability properties. Throughout the paper, we shall assume for simplicity that the manifolds, metrics and maps being considered are smooth, although many of the results that we quote or obtain can be shown to hold under weaker differentiability properties. Recall \cite{BCR1-2002, Eis1934, KaMi1980, Sta1897} that a {\emph {St\"ackel metric}} on an $n$-dimensional manifold $M$ is a Riemannian metric $g$ for which there exist local coordinates $\big(x^{1},\dots,x^{n}\big)$ in which the metric has the expression
\begin{gather}\label{Stackelform}
{\rm d}s^{2}=g_{ij}{\rm d}x^{i}{\rm d}x^{j}=\frac{\det S}{s^{11}}\big({\rm d}x^{1}\big)^{2}+\cdots+\frac{\det S}{s^{n1}}\big({\rm d}x^{n}\big)^{2},
\end{gather}
where $S$ is a {\emph{St\"ackel matrix}}, that is a non-singular $n\times n$ matrix $S=(s_{ij})$ of the form
\begin{gather}\label{StackelMatrix}
S=\left(
\begin{matrix}
s_{11}\big(x^{1}\big)&\dots&s_{1n}\big(x^{1}\big) \\
\vdots& & \vdots \\
s_{n1}\big(x^{n}\big)&\dots&s_{nn}\big(x^{n}\big)
\end{matrix}
\right),
\end{gather}
and $s^{ij}$ denotes the cofactor of the component $s_{ij}$ of the matrix~$S$. St\"ackel matrices thus have the property that for each $1\leq i \leq n$, their $i$-th row depends only on the $i$-th local coordinate~$x^{i}$, and that the cofactor $s^{ij}$ is independent of the $i$-th local coordinate $x^{i}$. Furthermore, the diagonal components of the St\"ackel metric~(\ref{Stackelform}) are given by the inverses of the entries of the first row of the inverse St\"ackel matrix $A=S^{-1}$.

The importance of St\"ackel metrics stems from the fact that they constitute the most general class of Riemannian metrics admitting orthogonal coordinates for which the Hamilton--Jacobi equation
\begin{gather}\label{HJ}
g^{ij}\partial_{i}W\partial_{j}W=E,
\end{gather}
for the geodesic flow of $(M,g)$, where $E$ denotes a positive real constant, admits a complete integral obtained by additive separation of variables into ordinary differential equations. It is useful at this stage to recall that a complete integral of (\ref{HJ}) is defined as a parametrized family of solutions
\begin{gather}\label{complint}
W=W\big(x^{1},\dots,x^{n};a_{1},\dots,a_{n}\big),\qquad a_{1}:=E,
\end{gather}
defined over the domain $U\subset M$ of the local coordinates $\big(x^{1},\dots,x^{n}\big)$ and depending smoothly on $n$ parameters $(a_{1},\dots,a_{n})$ defined on an open subset $A\subset {\mathbb{R}}^{n}$, such that the rank condition
\begin{gather*}%\label{HJrank}
\det \left(\frac{\partial^{2}W}{\partial x^{i}\partial a_{j}}\right) \neq 0,
\end{gather*}
is satisfied throughout the open set $U\times A$.

It is easily verified that the Hamilton--Jacobi equation (\ref{HJ}) will admit an additively separable complete integral $W\big(x^{1},\dots, x^{n};a_{1},\dots a_{n}\big)$ of the form
\[
W=W_{1}\big(x^{1};a_{1},\dots, a_{n}\big)+\cdots+W_{n}\big(x^{n};a_{1},\dots, a_{n}\big),
\]
if and only if the summands $W_{\alpha}$ satisfy the set of separated ordinary differential equations given by
\[
\left(\frac{{\rm d}W_{i}}{{\rm d}x^{i}}\right)^{2}=\sum_{j=1}^{n}s_{ij}\big(x^{i}\big)a_{j}.
\]
The $n$ parameters $(a_{1},\dots , a_{n})$ appearing in the expression of the additively separable complete integral~(\ref{complint}) thus correspond to the separation constants arising from the complete separation of variables of the Hamilton--Jacobi equation into ordinary differential equations.
One of the key consequences of this complete separation of variables property is that the geodesic flow of an $n$-dimensional St\"ackel metric admits a linearly independent set of $n-1$ mutually Poisson-commuting quadratic first integrals, given by
\begin{gather}\label{FirstInt}
K_{(l)} = K_{(l)}^{ij} p_i p_j = \sum_{j=1}^{n}a_{lj} p_{j}^{2}, \qquad 2\leq l \leq n
\end{gather}
with
\[ \big\{K_{(l)},K_{(m)}\big\}=0,\qquad \text{for} \quad 1\leq l,m \leq n,
\]
where $A=(a_{ij})$ denotes as above the inverse of the St\"ackel matrix $S$ given by (\ref{StackelMatrix}). Note that with the notations of (\ref{FirstInt}), we have $K_{(1)}=H$, where
\[
H=g^{ij}p_{i}p_{j},
\]
denotes the Hamiltonian for the geodesic flow.

 A question closely related to the additive separation of the Hamilton--Jacobi equation is that of the complete multiplicative separation of the Helmholtz equation
 \begin{gather}\label{Helm}
-\Delta_{g} u =\lambda u,
\end{gather}
where
\[
\Delta_{g} = \frac{1}{\sqrt{|g|}}\partial_{i}\Big(\sqrt{|g|}g^{ij}\partial_{j}\Big),\qquad |g|=\det(g_{ij}),
\]
denotes the Laplace--Beltrami operator on~$(M,g)$ and $\lambda$ denotes a non vanishing real constant, into ordinary differential equations. By complete multiplicative separation, we mean, follo\-wing~\cite{BCR1-2002}, the existence of a parametrized family of solutions $u$ defined on a domain $U\subset M$ with local coordinates $\big(x^{1},\dots, x^{n}\big)$ of the form
\[
u\big(x^{1},\dots, x^{n};a_{1},\dots, a_{2n}\big)=\prod_{i=1}^{n}u_{i}\big(x^{1},\dots, x^{n};a_{1},\dots, a_{2n}\big),
\]
depending smoothly on $2n$ parameters $(a_{1},\dots, a_{2n})$ defined on open subset $A\subset {\mathbb R}^{2n}$, satisfying the rank condition
\begin{gather*}%\label{Helmrank}
\det \left(
\begin{matrix}
\dfrac{\partial v_{1}}{\partial a_{1}}&\dots&\dfrac{\partial v_{1}}{\partial a_{2n}} \\
\vdots& & \vdots \\
\dfrac{\partial v_{n}}{\partial a_{1}}&\dots&\dfrac {\partial v_{n}}{\partial a_{2n}}\vspace{1mm}\\
\dfrac{\partial w_{1}}{\partial a_{1}}&\dots&\dfrac{\partial w_{1}}{\partial a_{2n}} \\
\vdots& & \vdots \\
\dfrac{\partial w_{n}}{\partial a_{1}}&\dots&\dfrac {\partial w_{n}}{\partial a_{2n}}
\end{matrix}
\right) \neq 0,
\end{gather*}
at every point of $U\times A$, where
\[
v_{i}=\frac{u_{i}'}{u_{i}} \qquad w_{i}=\frac{u_{i}''}{u_{i}}.
\]
This separation requires that additional conditions, known as the {\emph{Robertson conditions}}, and given by
\[
\partial_{i}\Gamma_{j}=0,\qquad 1\leq i\neq j\leq n,
\]
where
\[
\Gamma_{i}=-\partial_{i}\left[\log \frac{(\det S)^{\frac{n}{2}-1}s^{i1}}{\big(s^{11}\cdots s^{n1}\big)^{\frac{1}{2}}}\right],
\]
be satisfied. We refer to \cite{BCR1-2002,BCR2-2002,Eis1934,KaMi1980,KaMi1981,Rob1928} for detailed proofs of the fact that the Robertson conditions are necessary and sufficient for complete multiplicative separation of the Helmholtz equation. The Robertson conditions were interpreted by Eisenhart \cite{Eis1934} in terms of the vanishing of the off-diagonal components of the Ricci tensor of the underlying St\"ackel metric, that is
\begin{gather}\label{classicalRobertson}
R_{ij}=0\qquad \text{for}\qquad 1\leq i\neq j\leq n.
\end{gather}
When the Robertson conditions are satisfied, the Poisson-commuting quadratic first integrals of the geodesic flow give rise to $n-1$ linearly independent second-order differential operators which commute with the Laplace--Beltrami operator and also commute pairwise. Rewriting the quadratic first integrals $K_{(l)}$ defined by (\ref{FirstInt}) in the form
\[
K_{(l)}=K_{(l)}^{ij}p_{i}p_{j},
\]
these commuting operators, denoted by $\Delta_{K_{(l)}}$, are of the form
\[
\Delta_{K_{(l)}}=\nabla_{i}\big(K_{(l)}^{ij}\nabla_{j}\big),\qquad 2\leq l \leq n,
\]
where $\nabla_i$ denotes the Levi-Civita connection on $(M,g)$. These operators, which are often referred to as {\emph{symmetry operators}}, admit the separable solutions of the Helmholtz equation as (formal) eigenfunctions. We will not give any further details on symmetry operators at this stage, nor shall we say anything about the proofs of the results we have just recalled since we shall shortly state and prove generalizations of these to the case of Painlev\'e metrics, which admit all St\"ackel metrics as a special case.

We conclude these preliminaries by remarking that the above setting may be expanded significantly by considering conformal deformations of St\"ackel metrics which are compatible with the complete separation of the Helmholtz equation into ordinary differential equations, thus giving rise to the more general notion of {\emph{$R$-separability}} for the Helmholtz equation. Again, we shall not give any additional details on these topics at this stage since conformal deformations and $R$-separability will be studied in the remainder of this paper in the more general setting of Painlev\'e metrics. We refer to \cite{BCR1-2002, BCR2-2002, BCR2005, CR2006,CR2007, KaMi1980, KaMi1981, KaMi1984} for lucid accounts of the key results on the separability and $R$-separability properties of St\"ackel metrics, their connection to Killing tensors, quadratic first integrals of the geodesic flow and symmetry operators for the Laplace--Beltrami operator. We also refer to \cite{Ben2016, Mi1988} for recent surveys on separability on Riemannian manifolds and to \cite{Car1977} for a penetrating analysis of the relations between quadratic first integrals of the geodesic flow, symmetry operators and conserved currents, in the general setting of Riemannian or pseudo-Riemannian manifolds.

With these preliminaries at hand, we are now ready to introduce the class of Painlev\'e met\-rics~\cite{Pain1897, Per1990}. As stated above, Painlev\'e metrics arise as a natural generalization of St\"ackel metrics to the case in which one no longer seeks complete separation of the Hamilton--Jacobi equation into ordinary differential equations, but rather separation into partial differential equations involving groups of variables. The separable coordinates admitted by Painlev\'e metrics are thus generally {\emph {not orthogonal}}, although they are orthogonal with respect to groups of variables. Let us recall that our goal in this paper is to carry out for Painlev\'e metrics the analogue of the separability and $R$-separability analyses of the Helmholtz equation which has been extensively worked out for St\"ackel metrics in \cite{BCR1-2002,BCR2-2002,CR2006,KaMi1980,KaMi1981, KaMi1984}, and to show that the separability into groups of variables gives rise to vector spaces of mutually commuting symmetry operators for the Laplace--Beltrami operator, the dimension of which is determined by the number of groups of variables. In particular, we will generalize to the case of Painlev\'e metrics the Robertson conditions and the characterization thereof in terms of the Ricci tensor. We now proceed to define the class of Painlev\'e metrics along lines similar to the ones used above for St\"ackel metrics.

Let $(M,g)$ be an $n$-dimensional Riemannian manifold and let ${\bf x}=\big(x^{1},\dots , x^{n}\big)$ denote a set of local coordinates on $M$. We shall consider partitions of ${\bf x}$ into $r\geq 2$ groups of local coordinates,
\[
{\bf x}=\big({{\bf x}^{1}},\dots,{{\bf x}^{r}}\big),
\]
where
\[
{\bf {x}}^{\alpha}=\big(x^{i_{\alpha}}\big),\qquad 1_{\alpha}\leq i_{\alpha} \leq l_{\alpha},\qquad \text{and}\qquad \sum_{\alpha=1}^{r}||l_{\alpha}||=n,
\]
where $||l_{\alpha}||$ denotes the number of local coordinates in the group of label $\alpha$. Latin indices $1\leq i,j,\ldots \leq n$ will be used to label the local coordinates on $M$, greek indices $\alpha,\beta,\dots$ to label the $r$ groups of local coordinates, and hybrid indices $i_{\alpha}$, $1_{\alpha} \leq i_{\alpha}\leq l_{\alpha}$ to denote the local coordinates within the group $\bf{x^{\alpha}}$. Unless there is an ambiguity in the notation being used, in which case we will write out the summation signs explicitly, we shall apply the summation convention with the above range of indices.
A {\em generalized St\"ackel matrix} is a non-singular $r\times r$ matrix-valued function $S$ on $M$ of the form
\begin{gather}\label{genStackelmat}
S=\left(
\begin{matrix}
s_{11}\big({{\bf x}^{1}}\big)&\dots&s_{1r}\big({{\bf x}^{1}}\big) \\
\vdots& & \vdots \\
s_{r1}\big({{\bf x}^{r}}\big)&\dots&s_{rr}\big({\bf x}^{r}\big)
\end{matrix}
\right),
\end{gather}
where for each $1\leq \alpha \leq r$,
\[
{\bf x}^{\alpha}=\big(x^{i_{\alpha}}\big),\qquad 1_{\alpha}\leq i_{\alpha} \leq l_{\alpha}.
\]
Let $s^{\alpha \beta}$ denote the cofactor of the component $s_{\alpha \beta}$ of $S$. We note that the cofactor $s^{\beta \gamma}$ will not depend on the group of variables ${\bf x}^{\beta}=\big(x^{i_{\beta}}\big)$, $1_{\beta}\leq i_{\beta} \leq l_{\beta}$. Moreover we shall assume that
\begin{gather} \label{StackelPositiveness}
\frac{\det S}{s^{\alpha 1}} > 0, \qquad \forall\, 1\leq \alpha \leq r ,
\end{gather}
in order to work with {\emph{Riemannian}} Painlev\'e metrics.

\begin{Definition} \label{Painleve}
Let $S$ be a generalized St\"ackel matrix satisfying (\ref{StackelPositiveness}). A Painlev\'e metric is a~Riemannian metric~$g$ for which there exist local coordinates ${\bf x}=\big({\bf x}^{1},\dots,{\bf x}^{r}\big)$ such that
\begin{gather}\label{Painleveform}
{\rm d}s^{2}=g_{ij}{\rm d}x^{i}{\rm d}x^{j}=\frac{\det S}{s^{11}} G_1 +\cdots+\frac{\det S}{s^{r1}} G_r,
\end{gather}
where each of the quadratic differential forms
\begin{gather}\label{blockmetric}
G_{\alpha}=G_{\alpha}({{\bf x}^{\alpha}})=\sum_{i_{\alpha}=1_{\alpha}}^{l_{\alpha}}\sum_{j_{\alpha}=1_{\alpha}}^{l_{\alpha}}G_{\alpha}\big({{\bf x}^{\alpha}}\big)_{i_{\alpha}j_{\alpha}}{\rm d}x^{i_{\alpha}}{\rm d}x^{j_{\alpha}},\qquad 1\leq \alpha \leq r,
\end{gather}
is positive-definite in its arguments and depends only on the group of variables ${\bf x}^{\alpha}$.
\end{Definition}
We may thus write the metric (\ref{Painleveform}) in block-diagonal form as
\[
{\rm d}s^{2}=\sum_{\alpha =1}^{r}\sum_{\ialpha = 1_{\alpha}}^{\lalpha}\sum_{\jalpha = 1_{\alpha}}^{\lalpha}(g_{\alpha})_{\ialpha \jalpha}{\rm d}x^{\ialpha}{\rm d}x^{\jalpha}=\sum_{\alpha =1}^{r} \frac{\det S}{s^{\alpha 1}} \sum_{\ialpha = 1_{\alpha}}^{\lalpha}\sum_{\jalpha = 1_{\alpha}}^{\lalpha}(G_{\alpha})_{i_{\alpha}j_{\alpha}}{\rm d}x^{i_{\alpha}}{\rm d}x^{j_{\alpha}}.
\]
It is important to note that even though Painlev\'e metrics (\ref{Painleveform}) are block-diagonal, and each quadratic differential form (\ref{blockmetric}) defines a Riemannian metric on the submanifolds defined by the level sets ${\bf x}^{\beta}={\bf{c}}^{\beta}$, $\beta\neq \alpha$, {\emph {a Painlev\'e metric is generally not a direct sum of Riemannian metrics, nor a warped product, except for special non-generic cases}}. We also note that Painlev\'e metrics of semi-Riemannian (and in particular Lorentzian) signature can readily be defined by modifying the requirement that each of the quadratic differential forms $G_{\alpha}$ given by (\ref{blockmetric}) be positive-definite to one in which $G_{{\alpha}}$ is assumed to be of signature $(p_{\alpha},q_{\alpha})$ with $p_{\alpha}+q_{\alpha}=l_{\alpha}$. Finally we also remark that the Painlev\'e form (\ref{Painleveform}) is obviously invariant under smooth and invertible changes of coordinates of the form ${\tilde{\bf x}}^{\alpha}={\bf{f}}^{\alpha}({\bf x}^{\alpha})$, where $1\leq \alpha \leq r$.

Let us call \emph{block orthogonal coordinates} a system of coordinates $({\bf x}^{\alpha})$ such that the metric $g$ has the form
\begin{gather*} %\label{BlockOrthoCoord}
 g=\sum_{\alpha =1}^{r} c_\alpha G_\alpha = \sum_{\alpha =1}^{r} c_\alpha \sum_{\ialpha = 1_{\alpha}}^{\lalpha}\sum_{\jalpha = 1_{\alpha}}^{\lalpha}(G_{\alpha})_{i_{\alpha}j_{\alpha}}{\rm d}x^{i_{\alpha}}{\rm d}x^{j_{\alpha}},
\end{gather*}
where $c_\alpha$ are non-vanishing scalar functions on $M$ and the metrics $G_\alpha$ are given by (\ref{blockmetric}).
In analogy with the St\"ackel case, we have

\begin{Proposition} \label{HJSeparability-Painleve}A metric $g$ is of the Painlev\'e form \eqref{Painleveform} if and only if there exist block orthogonal coordinates such that the Hamilton--Jacobi equation
\[
g^{ij}\partial_{i}W\partial_{j}W=E,
\]
admits a parametrized family of solutions which is sum-separable into groups of variables, of the form
\begin{gather}\label{Painsepform}
W=W_{1}\big({\bf {x}}^{1};{{a}}_{1},\dots ,{{a}}_{r}\big)+\cdots+W_{r}\big({\bf {x}}^{r};{{{a}}_{1},\dots, {a}}_{r}\big),
\end{gather}
depending smoothly on $r$ arbitrary real constants $({{a}}_{1}:=E,a_{2},\dots ,{{a}}_{r})$ defined on an open subset $A\subset \mathbb{R}^{r}$, and satisfying the rank condition
\begin{gather}\label{rankPainHJ}
\det \left(\sum_{\ialpha =1_{\alpha}}^{\lalpha}\sum_{\jalpha =1_{\alpha}}^{\lalpha}\big(G^{\alpha}\big)^{\ialpha \jalpha}\left(\frac{\partial W_{\alpha}}{\partial x^{\ialpha}}\right)\left( \frac{\partial^{2}W_{\alpha}}{\partial a_{\gamma}\partial x^{\jalpha}}\right)\right)\neq 0,
\end{gather}
where
\[
G^{\alpha} = (G_{\alpha})^{-1 }.
\]
\end{Proposition}
This Proposition will be proved in Section \ref{PainHJSection} as well as other (intrinsic) characterizations of Painlev\'e metrics.

In further analogy with the St\"ackel case, we now recall that Painlev\'e metrics admit $r$ linearly independent quadratic first integrals of the geodesic flow which are Poisson commuting. Indeed, the summands $W_{\alpha}$ appearing in (\ref{Painsepform}) satisfy the following set of first-order PDEs \cite{Per1990}
\begin{gather}
{\mathcal{F}}_{1}\left({\bf {x}}^{1},\frac{\partial W_{1}}{\partial {\bf x}^{1}}\right) =a_{1}s_{11}\big({\bf {x}}^{1}\big)+\cdots +a_{r}s_{r1}\big({\bf {x}}^{1}\big),\nonumber\\ \quad \vdots \nonumber\\ {\mathcal{F}}_{r}\left({\bf {x}}^{r},\frac{\partial W_{r}}{\partial {\bf x}^{r}}\right) =a_{1}s_{r1}\big({\bf {x}}^{r}\big)+\cdots +a_{r}s_{rr}\big({\bf {x}}^{r}\big),\label{HJsep}
\end{gather}
where
\begin{gather}\label{F}
{\mathcal{F}}_{\alpha}=\big(G^{\alpha}\big)^{i_{\alpha}j_{\alpha}}\big({\bf {x}}^{\alpha}\big)p_{i_{\alpha}}p_{j_{\alpha}},\qquad 1\leq \alpha \leq r,
\end{gather}
and where the $(a_{\alpha})$ are arbitrary real separation constants.
It follows now directly from the separated equations (\ref{HJsep}) and from the fact that the generalized St\"ackel matrix $S$ is non-singular that
one obtains $r$ linearly independent Poisson-commuting quadratic first integrals~$K_{(\alpha)}$ of the geodesic flow by solving for the $r$ separation constants $(a_{\alpha})$ from the separated equations (\ref{HJsep}). These are explicitly given by (see~\cite{Per1990})
\[
K_{(\alpha)}=\sum_{\beta =1}^{r}\big(S^{-1}\big)_{\alpha \beta}\mathcal{F}_{\beta},\qquad 1\leq \beta \leq r,
\]
or equivalently
\[
K_{(\alpha)}=K_{(\alpha)}^{ij}p_{i}p_{j},
\]
where for each $1\leq \alpha \leq r$, $\big(K_{(\alpha)}^{ij}\big)$ is a symmetric block-diagonal tensor defined by
\begin{gather}\label{Killingtensor}
K_{(\alpha)}^{i_{\beta}j_{\beta}}=\frac{s^{\beta \alpha}}{\det S}\big(G^{\beta}\big)^{i_{\beta}j_{\beta}},\qquad K_{(\alpha)}^{i_{\beta}j_{\gamma}}=0\qquad \text{for all}\qquad 1 \leq \beta \neq \gamma \leq r.
\end{gather}
These quadratic first integrals satisfy
\begin{gather}\label{PoissonCommute}
\big\{K_{(\alpha)}, K_{(\beta)}\big\}=0, \qquad 1\leq \alpha,\beta \leq r, \qquad \text{where}\qquad K_{(1)}=H.
\end{gather}
The condition $\{K_{(\alpha)},H\}=0$ is equivalent to $\big(K_{(\alpha)}^{ij}\big)$ being a symmetric {\emph{Killing tensor}}, that is
\[
\nabla_{i}K_{(\alpha) jk}+\nabla_{j}K_{(\alpha) ki}+\nabla_{k}K_{(\alpha) ij}=0.
\]
The commutation relations (\ref{PoissonCommute}) are thus equivalent to the vanishing of the Schouten brackets of the pairs of Killing tensors $(K_{(\alpha) ij}),(K_{(\beta) ij})$.

There exist a few classical examples of Painlev\'e metrics in the litterature. They include for instance the {\emph{di Pirro metrics}} \cite{diPirro1896, Per1990}, for which the Hamiltonian of the geodesic flow is of the form
\begin{gather}\label{diPirro}
H=g^{ij}p_{i}p_{j}=\big(c_{12}\big(x^{1},x^{2}\big)+c_{3}\big(x^{3}\big)\big)^{-1}\big[a_{1}\big(x^{1},x^{2}\big)p_{1}^{2}+a_{2}\big(x^{1},x^{2}\big)p_{2}^{2} +a_{3}\big(x^{3}\big)p_{3}^{2}\big].
\end{gather}
It may indeed be verified directly that the function
\[
K=\big(c_{12}\big(x^{1},x^{2}\big)+c_{3}\big(x^{3}\big)\big)^{-1}\big[c_{3}\big(x^{3}\big)(a_{1}\big(x^{1},x^{2}\big)p_{1}^{2} +a_{2}\big(x^{1},x^{2}\big)p_{2}^{2})-c_{12}\big(x^{1},x^{2}\big)a_{3}\big(x^{3}\big)p_{3}^{2}\big].
\]
Poisson-commutes with $H$, and thus defines a Killing tensor, which together with the metric tensor generates a maximal linearly independent set of Killing tensors for generic choices of the metric functions $c_{12}$, $a_{1}$, $a_{2}$, $a_{3}$, $c_{3}$ in~(\ref{diPirro}). Painlev\'e metrics also appear in the context of geodesically equivalent metrics as metrics admitting projective symmetries, see \cite{Top1999, Top2002}, and also as instances of 4-dimensional Lorentzian metrics admitting a Killing tensor \cite{HaMa1978} (see Section~\ref{Perspectives} for further remarks on the latter point). Finally, we mention the recent paper by Chanu and Rastelli~\cite{CR2019} that provides a classification of Painlev\'e metrics with vanishing Riemann tensor in dimension~$3$,  i.e., in~$\mathbb{E}_3$. We will give some examples of Painlev\'e metrics in all dimensions satisfying the generalized Robertson conditions (see below) as well as a catalogue of such metrics in dimensions $2$, $3$, $4$ in Section \ref{2}.

As we stated above, our main goal in this paper is to investigate for the class of Painlev\'e metrics the closely related question of product separability for the Helmholtz equation (\ref{Helm}), and the relationship between quadratic first integrals of the geodesic flow and symmetry operators for the Laplace--Beltrami equation. The Laplace--Beltrami operator for a Painlev\'e metric $g$ given by~(\ref{Painleveform}) can be expressed in terms of the generalized St\"ackel matrix $S$ and the Laplace--Beltrami operators for the $r$ Riemannian metrics $G_{\beta}$, $1\leq \beta \leq r$, defined by~(\ref{blockmetric}), corresponding to the blocks of variables ${\bf x}^{\beta}$, $1\leq \beta \leq r$. We have
\begin{gather}
\Delta_{g} u =\sum_{\beta=1}^{r}\Bigg[\left(\frac{s^{\beta 1}}{\det S}\right)\Bigg\{\Delta_{G_{\beta}}u\nonumber\\
\hphantom{\Delta_{g} u =}{}+\sum_{i_{\beta}=1_{\beta}}^{l_{\beta}}\sum_{{j}_{\beta}=1_{\beta}}^{l_{\beta}}\big(G^{\beta}\big)^{i_{\beta} j_{\beta}}\left[\partial_{i_{\beta}}\left(\log \frac{(\det S)^{\frac{n}{2}-1}s^{\beta 1}}{\big(s^{11}\big)^{\frac{l_1}{2}} \cdots \big(s^{r1}\big)^{\frac{l_r}{2}}}\right)\right]\partial_{j_{\beta}}u\Bigg\}\Bigg],\label{Laplacian}
\end{gather}
where $\Delta_{G_{\beta}}$ denotes the Laplace--Beltrami operator for the Riemannian metric $G_{\beta}$, that is
\begin{gather*}%\label{blockLaplacian}
\Delta_{G_{\beta}} = \sum_{i_{\beta}=1}^{\lbeta}\sum_{j_{\beta}=1}^{\lbeta}\frac{1}{\sqrt{|G_{\beta}|}}\partial_{i_{\beta}}\Big(\sqrt{|G_{\beta}|}
\big(G^{\beta}\big)^{i_{\beta}j_{\beta}}\partial_{j_{\beta}}\Big),\qquad |G_{\beta}|=\det((G_{\beta})_{i_{\beta}j_{\beta}}).
\end{gather*}

We now state our main results, the proofs of which will be given in Section \ref{ProofsofMainTheorems}. We first define the \emph{generalized Robertson conditions}, in analogy with the classical Robertson conditions (\ref{classicalRobertson}) for St\"ackel metrics.

\begin{Definition} \label{defiRC}
A Painlev\'e metric $g$ is said to satisfy the generalized Robertson conditions if and only if the differential conditions
	\begin{gather}\label{GenRobertson}
\partial_{j_{\alpha}}\gamma_{i_{\beta}}=0,\quad \textrm{for all} \quad 1\leq \alpha \neq \beta \leq r, \quad 1_{\alpha}\leq i_{\alpha} \leq l_{\alpha},\quad 1_{\beta}\leq i_{\beta} \leq l_{\beta},
\end{gather}
where
\[
\gamma_{i_{\beta}}=-\partial_{i_{\beta}}\left[\log \frac{(\det S)^{\frac{n}{2}-1}s^{\beta 1}}{\big(s^{11}\big)^{\frac{l_1}{2}} \cdots \big(s^{r1}\big)^{\frac{l_r}{2}}}\right],
\]
are satisfied.
\end{Definition}
The generalized Robertson conditions (\ref{GenRobertson}) imply that
\begin{gather}\label{GenRobertsonsimpl}
\partial_{\ialpha}\gamma^{\jbeta}=0, \quad \textrm{for all} \quad 1\leq \alpha \neq \beta \leq r,\\
\label{gammaup}
\gamma^{j_{\beta}}:= \sum_{i_{\beta}=1_{\beta}}^{l_{\beta}}\big(G^{\beta}\big)^{i_{\beta}j_{\beta}} \gamma_{i_{\beta}}.
\end{gather}
We shall be working with both the forms (\ref{GenRobertson}) and (\ref{GenRobertsonsimpl}) of these conditions. Note that if the Robertson conditions hold, then the positive Laplace--Beltrami operator can be written in a synthetic form as
\begin{gather*}%\label{LaplacianSimp}
 -\Delta_{g} u  =   \sum_{\beta=1}^{r} \left(\frac{s^{\beta 1}}{\det S}\right)\left\{-\Delta_{G_{\beta}}u + \sum_{{j}_{\beta}=1_{\beta}}^{l_{\beta}} \gamma^{j_{\beta}} \partial_{j_{\beta}}u\right\}
   =   \sum_{\beta=1}^{r} \left(\frac{s^{\beta 1}}{\det S}\right) B_\beta u , \nonumber
\end{gather*}
where the operators $B_\beta = -\Delta_{G_{\beta}} + \sum\limits_{{j}_{\beta}=1_{\beta}}^{l_{\beta}} \gamma^{j_{\beta}} \partial_{j_{\beta}}$ only depend on the group of variables $\textbf{x}^\beta$.

As will be seen in Section \ref{PainHelmSection}, the generalization of the notion of complete multiplicative separation for the Helmholtz equation to the case of separation in terms of groups of variables is given by considering a parametrized family of product-separable solutions of the form
\begin{gather} \label{SeparabilitySchro}
u=\prod_{\beta =1}^{r}u_{\beta}\big({\bf x}^{\beta};a_{1},\dots,a_{r}\big),
\end{gather}
satisfying the rank condition
\begin{gather} \label{RankPainSchro}
 \det  \left( \partial_{a_\alpha} \left( \frac{B_\beta u_\beta}{u_\beta} \right) \right) \ne 0,
\end{gather}
where we assume that $u_{\beta}\neq 0$.

Our first result states the separability conditions for the Helmholtz equation, and gives their interpretation in terms of the vanishing of the off-block diagonal components of the Ricci tensor:

